% !TeX spellcheck = ru_RU
\documentclass[a4paper, 14pt, unknownkeysallowed]{extreport}
\usepackage[utf8]{inputenc}
\usepackage[T2A]{fontenc}
\usepackage[english,russian]{babel}
\usepackage{cmap}
\usepackage{enumitem}

\usepackage{csvsimple}
\usepackage{hyphenat} 
\usepackage{amstext, amsmath,amsfonts,amssymb,amsthm,mathtools} 

\usepackage{setspace}
\onehalfspacing 

\usepackage{geometry}
\geometry{left=30mm}
\geometry{right=10mm}
\geometry{top=10mm}
\geometry{bottom=20mm}

\usepackage{graphicx}
\graphicspath{{img/}} 

\usepackage{threeparttable}
\usepackage{bigdelim}

\usepackage{setspace}
\onehalfspacing % Полуторный интервал

\frenchspacing
\usepackage{indentfirst} % Красная строка

\usepackage{cases}

\usepackage[unicode,pdftex]{hyperref} % Ссылки в pdf
\hypersetup{hidelinks}

\usepackage{xcolor}
\usepackage{listings}
% Для листинга кода:
\lstset{%
	language=lisp,   					% выбор языка для подсветки	
	basicstyle=\small\sffamily,			% размер и начертание шрифта для подсветки кода
	numbers=left,						% где поставить нумерацию строк (слева\справа)
	%numberstyle=,					% размер шрифта для номеров строк
	stepnumber=1,						% размер шага между двумя номерами строк
	numbersep=5pt,						% как далеко отстоят номера строк от подсвечиваемого кода
	frame=single,						% рисовать рамку вокруг кода
	tabsize=4,							% размер табуляции по умолчанию равен 4 пробелам
	captionpos=t,						% позиция заголовка вверху [t] или внизу [b]
	breaklines=true,					
	breakatwhitespace=true,				% переносить строки только если есть пробел
	escapeinside={\#*}{*)},				% если нужно добавить комментарии в коде
	backgroundcolor=\color{white},
}
\usepackage{caption}
\captionsetup{labelsep = endash}
\captionsetup[figure]{name = {}, justification = centerlast,labelformat=empty}
\captionsetup[lstlisting]{justification = justified}
\captionsetup[table]{justification = justified}


\usepackage{titlesec}
\titleformat{\chapter}{\LARGE\bfseries}{\thechapter}{20pt}{\LARGE\bfseries}
\titleformat{\section}{\Large\bfseries}{\thesection}{20pt}{\Large\bfseries}

\renewcommand\labelitemi{\ --\ }
\renewcommand\labelenumi{\theenumi)}

\newcommand{\img}[3] {
	\begin{figure}[h!]
		\includegraphics[scale = #1]{img/#2}
		\caption{#3}
		\label{img:#2}
	\end{figure}
}

\newcommand{\lst}[4]{
\lstinputlisting[language=Lisp, firstline=#1, lastline=#2, label=lst:#3,caption=#4]{sx.lsp}
}

\usepackage{siunitx,array,booktabs}


\begin{document}
% !TeX spellcheck = ru_RU
\begin{titlepage}
	
	\newgeometry{pdftex, left=2cm, right=2cm, top=2.5cm, bottom=2.5cm}
	\fontsize{12pt}{12pt}\selectfont
	\noindent \begin{minipage}{0.15\textwidth}
		\includegraphics[width=\linewidth]{img/b_logo.jpg}
	\end{minipage}
	\noindent\begin{minipage}{0.9\textwidth}\centering
		\textbf{Министерство науки и высшего образования Российской Федерации}\\
		\textbf{Федеральное государственное бюджетное образовательное учреждение высшего образования}\\
		\textbf{«Московский государственный технический университет имени Н.~Э.~Баумана}\\
		\textbf{(национальный исследовательский университет)»}\\
		\textbf{(МГТУ им. Н.~Э.~Баумана)}
	\end{minipage}

	\noindent\rule{18cm}{3pt}
	\newline\newline
	\noindent ФАКУЛЬТЕТ $\underline{\text{«Информатика и системы управления»~~~~~~~~~~~~~~~~~~~~~~~~~~~~~~~~~~~~~~~~~~~~~~~~~~~~~~~}}$ \newline\newline
	\noindent КАФЕДРА $\underline{\text{«Программное обеспечение ЭВМ и информационные технологии»~~~~~~~~~~~~~~~~~~~~~~~}}$\newline\newline\newline\newline\newline\newline\newline
	
	
	\begin{center}
		\Large\textbf{Исследование быстродействия умножения}
		
		\Large\textbf{матриц на языке Common Lisp}
%		
		\textbf{\newline}		
	\end{center}
	
	\noindent\textbf{Студент} $\underline{\text{~~Виноградов А. О.~~~~~~~~~~~~~~~~~~~~~~~~~~~~~~~~~~~~~~~~~~~~~~~~~~~~~~~~~~~~~~~~~~~~~~~~~~~~~~~~~~~~~~}}$\newline\newline
	\noindent\textbf{Группа} $\underline{\text{~~ИУ7-66Б~~~~~~~~~~~~~~~~~~~~~~~~~~~~~~~~~~~~~~~~~~~~~~~~~~~~~~~~~~~~~~~~~~~~~~~~~~~~~~~~~~~~~~~~~~~~~~~~~~~~}}$\newline\newline
	\noindent\textbf{Преподаватель} $\underline{\text{~~Строганов Ю.В.~~~~~~~~~~~~~~~~~~~~~~~~~~~~~~~~~~~~~~~~~~~~~~~~~~~~~~~~~~~~~~~~~~~~~~~~~~~~~}}$\newline
	
	\begin{center}
		\vfill
		Москва~---~\the\year
		~г.
	\end{center}
	\restoregeometry
	
\end{titlepage}
\setcounter{page}{2}

\chapter*{Введение}
\addcontentsline{toc}{chapter}{Введение}
Цель работы: приобрести навыки работы в системе Common Lisp.

Задачи работы: изучить работу форм — функций, которые особым образом обрабатывают свои аргументы и особенности их работы в Lisp.

\chapter{Практические задания}

\section{Написать функцию, которая принимает целое число и возвращает первое четное число, не меньшее аргумента}
\begin{lstlisting}
(defun f1 (x)
	(+ x (mod x 2))
)
\end{lstlisting}

\section{Написать функцию, которая принимает число и возвращает число того же знака, но с модулем на 1 больше модуля аргумента}

\begin{lstlisting}
(defun f2 (x)
	(if (minusp x) (- x 1) (+ x 1))
)
\end{lstlisting}

\section{Написать функцию, которая принимает два числа и возвращает список из этих чисел, расположенный по возрастанию}

\begin{lstlisting}
(defun f3 (a b)
	(if (< a b) (list a b) (list b a))
)
\end{lstlisting}

\section{Написать функцию, которая принимает три числа и возвращает Т только тогда, когда первое число расположено между вторым и третьим}

\begin{lstlisting}
(defun f4 (a b c)
	(if 
		(or
			(and (< b a) (< a c))
			(and (> b a) (> a c))
		)
		T
		Nil
	)
)
\end{lstlisting}

\section{Каков результат вычисления следующих выражений}

\begin{lstlisting}
(and 'fee 'fie 'foe) ;foe
(or 'fee 'fie 'foe) ;fee
(or nil 'fie 'foe) ;fie
(and nil 'fie 'foe) ;nil
(and (equal 'abc 'abc) 'yes) ;yes
(or (equal 'abc 'abc) 'yes) ;T
\end{lstlisting}


\section{Написать предикат, который принимает два числа-аргумента и возвращает Т, если первое число не меньше второго}.

\begin{lstlisting}
(defun f6 (a b) 
	(>= a b)
)
\end{lstlisting} 

\section{Какой из следующих двух вариантов предиката ошибочен и почему?}

\begin{lstlisting}
(defun pred1 (x)
	(and (numberp x) (plusp x)))

(defun pred2 (x)
	(and (plusp x)(numberp x)))
\end{lstlisting} 

Ошибочен второй вариант предиката, так как в ходе его вычисления сначала будет вычислено (plusp x), что может привести к ошибке, если x не численного типа.

\section{Решить задачу 4, используя для ее решения конструкции: только IF, только COND, только AND/OR.}
Написать функцию, которая принимает три числа и возвращает Т только тогда, когда первое число расположено между вторым и третьим.
\begin{lstlisting}
(defun f81 (a b c )
	(if (< b a)
		(if (< a c)
		T
		Nil)
		Nil))
\end{lstlisting} 

\begin{lstlisting}
(defun f82 (a b c )
	(cond ((< b a) (cond ((< a c) T) (T Nil))) 
			(T Nil))
)
\end{lstlisting} 

\begin{lstlisting}
(defun f83 (a b c )
	(and (< b a) (< a c)))
\end{lstlisting} 

\section{Переписать функцию how-alike, приведенную в лекции и использующую COND, используя только конструкции IF, AND/OR.}

\begin{lstlisting}
(defun how_alike (x y)
	(cond ((or (= x y) (equal x y)) 'the_same)
		((and (oddp x) (oddp y)) 'both_odd)
		((and (evenp x) (evenp y)) 'both_even)
		(T 'difference)
))
\end{lstlisting} 

\begin{lstlisting}
(defun how_alike (x y)
	(if (if (= x y)
		(equal x y))
		'the_same
		(if (if (oddp x)
			(oddp y))
		'both_odd
		(if (if (evenp x)
			(evenp y))
		'both_even
		'difference
))))
\end{lstlisting} 

\begin{lstlisting}
(defun how_alike (x y)
	(or
		(and (= x y) (equal x y) 'the_same)
		(and (oddp x) (oddp y) 'both_odd)
		(and (evenp x) (evenp y) 'both_even)
		'difference
))
\end{lstlisting} 

\end{document}